% PUC-Rio 1999 — 19 questões (14 MCQ + 5 discursivas)
% Fonte: HTML de banco de questões externo
% Notas de transcrição:
%   - Q8: expoente "9²¦" no HTML original — inferido como 9^{25} pela resposta D=24
%   - Q16: "(1,04)^{1}²" no HTML — inferido como (1,04)^{12} pelo contexto e resposta

---
tipo: Múltipla Escolha
dificuldade: Média
nivel: médio
disciplina: Matemática
assunto: Radiciação
gabarito: B
tags: [dízima periódica, simplificação]
fonte: concurso
concurso: PUC
banca: PUC
ano: 1999
numero: 1
---
\question
O valor de $\dfrac{\sqrt{1{,}\overline{7}}}{\sqrt{0{,}\overline{1}}}$ é:
\begin{choices}
\choice $4{,}\overline{4}$
\CorrectChoice $4$
\choice $4{,}\overline{7}$
\choice $3$
\choice $\dfrac{4}{3}$
\end{choices}

---
tipo: Múltipla Escolha
dificuldade: Média
nivel: médio
disciplina: Matemática
assunto: Radiciação
gabarito: A
tags: [raiz cúbica, raiz quadrada, potenciação]
fonte: concurso
concurso: PUC
banca: PUC
ano: 1999
numero: 2
---
\question
$\sqrt[3]{-8} \cdot \sqrt[2]{(-5)^2} =$
\begin{choices}
\CorrectChoice $-10$
\choice $-\sqrt{40}$
\choice $40$
\choice $\sqrt{40}$
\choice $2\sqrt{5}$
\end{choices}

---
tipo: Múltipla Escolha
dificuldade: Fácil
nivel: médio
disciplina: Matemática
assunto: Circunferência e Círculo
gabarito: C
tags: [área, comprimento, circunferência, raio]
fonte: concurso
concurso: PUC
banca: PUC
ano: 1999
numero: 3
---
\question
Triplicando-se o raio de uma circunferência:
\begin{choices}
\choice a área é multiplicada por $9\pi$.
\choice o comprimento é multiplicado por $3\pi$.
\CorrectChoice a área é multiplicada por $9$ e o comprimento por $3$.
\choice a área e o comprimento são ambos multiplicados por $3$.
\choice a área é multiplicada por $3$ e o comprimento por $9$.
\end{choices}

---
tipo: Múltipla Escolha
dificuldade: Média
nivel: médio
disciplina: Matemática
assunto: Combinações
gabarito: B
tags: [torneio, combinações, contagem]
fonte: concurso
concurso: PUC
banca: PUC
ano: 1999
numero: 4
---
\question
Um torneio de xadrez no qual cada jogador joga com todos os outros tem 351 partidas. O número de jogadores disputando é:
\begin{choices}
\choice $22$
\CorrectChoice $27$
\choice $26$
\choice $19$
\choice $23$
\end{choices}

---
tipo: Múltipla Escolha
dificuldade: Média
nivel: médio
disciplina: Matemática
assunto: Radiciação
gabarito: A
tags: [comparação de radicais, ordenação, números reais]
fonte: concurso
concurso: PUC
banca: PUC
ano: 1999
numero: 5
---
\question
Seja $a = 12(\sqrt{2}-1)$, $b = 4\sqrt{2}$ e $c = 3\sqrt{3}$. Então:
\begin{choices}
\CorrectChoice $a < c < b$
\choice $c < a < b$
\choice $a < b < c$
\choice $b < c < a$
\choice $b < a < c$
\end{choices}

---
tipo: Múltipla Escolha
dificuldade: Média
nivel: médio
disciplina: Matemática
assunto: Polinômios
gabarito: D
tags: [teorema do resto, divisão de polinômios]
fonte: concurso
concurso: PUC
banca: PUC
ano: 1999
numero: 6
---
\question
O resto da divisão do polinômio $x^3 + px + q$ por $x + 1$ é $4$ e o resto da divisão deste mesmo polinômio por $x - 1$ é $8$. O valor de $p$ é:
\begin{choices}
\choice $5$
\choice $-4$
\choice $0$
\CorrectChoice $1$
\choice $8$
\end{choices}

---
tipo: Múltipla Escolha
dificuldade: Média
nivel: médio
disciplina: Matemática
assunto: Teorema de Pitágoras
gabarito: C
tags: [triângulo retângulo, hipotenusa, sistema de equações]
fonte: concurso
concurso: PUC
banca: PUC
ano: 1999
numero: 7
---
\question
A hipotenusa de um triângulo retângulo mede $2\sqrt{61}$. A diferença entre os comprimentos dos dois outros lados é $2$. Então o menor lado tem comprimento:
\begin{choices}
\choice $\sqrt{30}$
\choice $7$
\CorrectChoice $10$
\choice $5\sqrt{6}$
\choice $11$
\end{choices}

---
tipo: Múltipla Escolha
dificuldade: Média
nivel: médio
disciplina: Matemática
assunto: Função Logarítmica
gabarito: D
tags: [número de algarismos, logaritmo, potência]
fonte: concurso
concurso: PUC
banca: PUC
ano: 1999
numero: 8
---
\question
Sabendo-se que $\log_{10} 3 \approx 0{,}47712$, podemos afirmar que o número de algarismos de $9^{25}$ é:
\begin{choices}
\choice $21$
\choice $22$
\choice $23$
\CorrectChoice $24$
\choice $25$
\end{choices}

---
tipo: Múltipla Escolha
dificuldade: Média
nivel: médio
disciplina: Matemática
assunto: Segmentos Proporcionais
gabarito: B
tags: [paralelogramo, ponto médio, razão, semelhança]
fonte: concurso
concurso: PUC
banca: PUC
ano: 1999
numero: 9
---
\question
$ABCD$ é um paralelogramo, $M$ é o ponto médio do lado $CD$, e $T$ é o ponto de interseção de $AM$ com $BD$. O valor da razão $\dfrac{DT}{BD}$ é:
\begin{choices}
\choice $\dfrac{1}{2}$
\CorrectChoice $\dfrac{1}{3}$
\choice $\dfrac{2}{5}$
\choice $\dfrac{1}{4}$
\choice $\dfrac{2}{7}$
\end{choices}

---
tipo: Múltipla Escolha
dificuldade: Fácil
nivel: médio
disciplina: Matemática
assunto: Função Quadrática
gabarito: C
tags: [parábola, interseção, sistemas]
fonte: concurso
concurso: PUC
banca: PUC
ano: 1999
numero: 10
---
\question
O número de pontos de interseção das duas parábolas $y = x^2$ e $y = 2x^2 - 1$ é:
\begin{choices}
\choice $0$
\choice $1$
\CorrectChoice $2$
\choice $3$
\choice $4$
\end{choices}

---
tipo: Múltipla Escolha
dificuldade: Média
nivel: médio
disciplina: Matemática
assunto: Área das Figuras Planas
gabarito: C
tags: [paralelogramo, área máxima, otimização]
fonte: concurso
concurso: PUC
banca: PUC
ano: 1999
numero: 11
---
\question
A área máxima de um paralelogramo com lados $a$, $b$, $a$, $b$ é:
\begin{choices}
\choice $a^2 + b^2$
\choice $2ab$
\CorrectChoice $ab$
\choice $a + b$
\choice $\dfrac{a}{b}$
\end{choices}

---
tipo: Múltipla Escolha
dificuldade: Média
nivel: médio
disciplina: Matemática
assunto: Reta
gabarito: D
tags: [colinearidade, determinante, geometria analítica]
fonte: concurso
concurso: PUC
banca: PUC
ano: 1999
numero: 12
---
\question
O valor de $x$ para que os pontos $(1,\,3)$, $(-2,\,4)$ e $(x,\,0)$ do plano sejam colineares é:
\begin{choices}
\choice $8$
\choice $9$
\choice $11$
\CorrectChoice $10$
\choice $5$
\end{choices}

---
tipo: Múltipla Escolha
dificuldade: Média
nivel: médio
disciplina: Matemática
assunto: Reta
gabarito: E
tags: [interseção de retas, equação da reta, geometria analítica]
fonte: concurso
concurso: PUC
banca: PUC
ano: 1999
numero: 13
---
\question
O ponto de interseção entre a reta que passa por $(4,\,4)$ e $(2,\,5)$ e a reta que passa por $(2,\,7)$ e $(4,\,3)$ é:
\begin{choices}
\choice $(3,\;5)$
\choice $(4,\;4)$
\choice $(3,\;4)$
\choice $\left(\dfrac{7}{2},\;4\right)$
\CorrectChoice $\left(\dfrac{10}{3},\;\dfrac{13}{3}\right)$
\end{choices}

---
tipo: Múltipla Escolha
dificuldade: Fácil
nivel: médio
disciplina: Matemática
assunto: Prismas
gabarito: A
tags: [paralelepípedo, diagonal, distância entre vértices]
fonte: concurso
concurso: PUC
banca: PUC
ano: 1999
numero: 14
---
\question
Considere um paralelepípedo retangular com lados $2$, $3$ e $6$ cm. A distância máxima entre dois vértices deste paralelepípedo é:
\begin{choices}
\CorrectChoice $7$ cm
\choice $8$ cm
\choice $9$ cm
\choice $10$ cm
\choice $11$ cm
\end{choices}

---
tipo: Discursiva
dificuldade: Média
nivel: médio
disciplina: Matemática
assunto: Pontos Notáveis do Triângulo
resposta: $AO = \dfrac{\sqrt{3}}{3}$ cm
tags: [triângulo equilátero, ortocentro, baricentro, altura]
fonte: concurso
concurso: PUC
banca: PUC
ano: 1999
numero: 15
---
\question
Seja $ABC$ um triângulo equilátero de lado $1$ cm em que $O$ é o ponto de encontro das alturas. Quanto mede o segmento $AO$?

---
tipo: Discursiva
dificuldade: Média
nivel: médio
disciplina: Matemática
assunto: Juros Compostos
resposta: A inflação acumulada será de $100 \cdot \left[(1{,}04)^{12} - 1\right]\%$.
tags: [inflação, juros compostos, montante]
fonte: concurso
concurso: PUC
banca: PUC
ano: 1999
numero: 16
---
\question
Suponha uma inflação mensal de $4\%$ durante um ano. De quanto será a inflação acumulada neste ano? (Pode deixar indicado o resultado.)

---
tipo: Discursiva
dificuldade: Fácil
nivel: médio
disciplina: Matemática
assunto: Equação do Segundo Grau
resposta: $a = -\dfrac{1}{4}$
tags: [raízes iguais, discriminante, delta]
fonte: concurso
concurso: PUC
banca: PUC
ano: 1999
numero: 17
---
\question
Quando o polinômio $x^2 + x - a$ tem raízes iguais?

---
tipo: Discursiva
dificuldade: Difícil
nivel: médio
disciplina: Matemática
assunto: Divisibilidade
resposta: Não é necessariamente verdadeiro. Contraexemplo: $a = 10$, pois $(100 + 10)^3 = 110^3 = 1\,331\,000$ é divisível por $100$, mas $100$ não divide $10$.
tags: [divisibilidade, contraexemplo, números naturais]
fonte: concurso
concurso: PUC
banca: PUC
ano: 1999
numero: 18
---
\question
Seja $a$ um número natural tal que $100$ é divisor de $(100+a)^3$. Então é necessariamente verdadeiro que $100$ é um divisor de $a$? Por quê?

---
tipo: Discursiva
dificuldade: Média
nivel: médio
disciplina: Matemática
assunto: Binômio de Newton
resposta: $8$
tags: [soma dos coeficientes, potência de polinômio, binômio]
fonte: concurso
concurso: PUC
banca: PUC
ano: 1999
numero: 19
---
\question
Ache a soma dos coeficientes do polinômio $(1 - 2x + 3x^2)^3$.
